% year: תשע"ח
% type: מבחן
% semester: א
% moed: א
% number: 3א
% points: 10
% categories: taylor, func-limits
% official_solution: yes
% source: exams/מבחנים/תשעח/מועד א תשעח.pdf

\begin{question}
(10 נק') חשבו את הגבול:
\[ \lim_{x\to0}\left(\frac{\sqrt{1-2x+x^3}-\sqrt[3]{1-3x+x^2}}{x^2}\right) \]
\underline{רמז}: אפשר להשתמש בנוסחאות מקלורין.
\end{question}

\begin{hint}
פתחו כל שורש לפי $(1+t)^p=1+pt+\frac{p(p-1)}{2}t^2+o(t^2)$ עד סדר $x^2$.
\end{hint}

\begin{solution}
נשתמש בפיתוח מקלורין $(1+t)^p=1+pt+\frac{p(p-1)}{2!}t^2+o(t^2)$ כאשר $t\to0$.

עבור השורש הראשון, $t=x^3-2x\to0$ ו-$p=\frac12$:
\begin{align*}
\sqrt{1-2x+x^3}&=1+\frac12(x^3-2x)-\frac18(x^3-2x)^2+o(x^2)\\
&=1-x+\frac{x^3}{2}-\frac18\left(4x^2-4x^4+x^6\right)+o(x^2)=1-x-\frac{x^2}{2}+o(x^2).
\end{align*}
(השתמשנו בכך ש-$t=O(x)$ ולכן $o(t^2)=o(x^2)$, ובכך שהחזקות $x^3,x^4,x^6$ הן $o(x^2)$.)

עבור השורש השני, $t=-3x+x^2$ ו-$p=\frac13$, $\frac{p(p-1)}{2}=\frac{\frac13\cdot(-\frac23)}{2}=-\frac19$:
\begin{align*}
\sqrt[3]{1-3x+x^2}&=1+\frac13(-3x+x^2)-\frac19(-3x+x^2)^2+o(x^2)\\
&=1-x+\frac{x^2}{3}-\frac19\left(9x^2-6x^3+x^4\right)+o(x^2)=1-x-\frac{2x^2}{3}+o(x^2).
\end{align*}
לכן
\[ \frac{\sqrt{1-2x+x^3}-\sqrt[3]{1-3x+x^2}}{x^2}=\frac{\left(-\frac12+\frac23\right)x^2+o(x^2)}{x^2}=\frac16+\frac{o(x^2)}{x^2}\xrightarrow[x\to0]{}\frac16. \]
\[ \boxed{\lim_{x\to0}\frac{\sqrt{1-2x+x^3}-\sqrt[3]{1-3x+x^2}}{x^2}=\frac16} \]
\end{solution}
